\section{ขอบเขตและข้อจำกัดของการวิจัย}

\subsection{ขอบเขตของการวิจัย (Scope)}

การวิจัยนี้มุ่งเน้นการพัฒนา ระบบวางแผนการท่องเที่ยวอัจฉริยะ (Smart Trip Planner) สำหรับจังหวัดขอนแก่น โดยประยุกต์ใช้เทคโนโลยี Large Language Models (LLMs) ร่วมกับสถาปัตยกรรม Retrieval-Augmented Generation (RAG) เพื่อสร้างแผนการเดินทางที่สอดคล้องกับความต้องการของผู้ใช้ และสามารถอ้างอิงข้อมูลจากฐานความรู้ที่ได้รับการจัดเก็บและปรับปรุงอย่างต่อเนื่อง

ระบบที่พัฒนาประกอบด้วยองค์ประกอบหลัก ดังนี้

\begin{enumerate}[leftmargin=1.8cm]
\item ระบบวางแผนการเดินทาง (Trip Planning System) ใช้เทคโนโลยี LLMs ร่วมกับ RAG ในการสร้างแผนการเดินทาง โดยพิจารณาจากข้อมูลที่ผู้ใช้กำหนด ได้แก่ วันที่เดินทาง รูปแบบการท่องเที่ยว (เช่น จำนวนสถานที่ที่ต้องการเข้าชม) ความสนใจ (เช่น วัด พิพิธภัณฑ์ ธรรมชาติ หรือร้านอาหาร) งบประมาณในการเดินทาง และสถานที่ที่ผู้ใช้ต้องการไปเป็นพิเศษ เพื่อสร้างแผนการเดินทางที่เหมาะสมและสอดคล้องกับข้อกำหนดของผู้ใช้
\item ระบบฐานความรู้และการสืบค้นข้อมูล (Knowledge Retrieval System) ใช้ฐานข้อมูลธรรมดา ร่วมกับฐานข้อมูลเวกเตอร์ (Vector Database) สำหรับจัดเก็บข้อมูลแหล่งท่องเที่ยว ร้านอาหาร กิจกรรม และเหตุการณ์ (Events) ภายในจังหวัดขอนแก่น โดยข้อมูลรวบรวมจากแหล่งข้อมูลสาธารณะ ได้แก่ Google Maps สื่อสังคมออนไลน์ และข้อมูลที่ผู้ดูแลระบบเพิ่มเข้าสู่ระบบ เพื่อสนับสนุนการค้นคืนข้อมูลสำหรับระบบ RAG และระบบตอบคำถามอัตโนมัติ
\item ระบบแชตบอต (Chatbot System) ใช้เทคโนโลยี LLMs และ RAG ในการตอบคำถามเกี่ยวกับการท่องเที่ยวจังหวัดขอนแก่น โดยอ้างอิงข้อมูลจากฐานความรู้ของระบบ เพื่อให้ข้อมูลที่เกี่ยวข้องกับสถานที่ท่องเที่ยว ร้านอาหาร กิจกรรม การเดินทาง และข้อมูลที่ผู้ดูแลระบบเผยแพร่
\end{enumerate}

เว็บแอปพลิเคชัน (Web-based Application) รองรับการใช้งานของผู้ใช้ 2 กลุ่ม ได้แก่

\begin{enumerate}[leftmargin=1.8cm]
\item ผู้ใช้งานทั่วไป สามารถค้นหาข้อมูลแหล่งท่องเที่ยว กิจกรรม และสร้างแผนการเดินทางผ่านเว็บไซต์ รวมถึงสามารถสมัครสมาชิกและเข้าสู่ระบบเพื่อใช้งานระบบสะสมคะแนนจากการสแกน QR Code ณ สถานที่ที่เข้าร่วมโครงการ ซึ่งเป็นความสามารถเพิ่มเติม (Optional Feature)
\item ผู้ดูแลระบบ (Administrator) สามารถจัดการข้อมูลแหล่งท่องเที่ยว ร้านอาหาร กิจกรรมหรืออีเวนต์ในจังหวัดขอนแก่น รวมถึงจัดการฐานความรู้ที่ใช้สำหรับระบบแชตบอต
\end{enumerate}

ขอบเขตของข้อมูลในการวิจัยครอบคลุม จังหวัดขอนแก่น โดยรวบรวมข้อมูลสถานที่ท่องเที่ยว ร้านอาหาร ที่พัก กิจกรรม และอีเวนต์ที่เกี่ยวข้องจากแหล่งข้อมูลสาธารณะและข้อมูลที่ผู้ดูแลระบบเพิ่มเข้าสู่ระบบ ทั้งนี้ การศึกษาให้ความสำคัญกับการรองรับการท่องเที่ยวทั้งในเขตอำเภอเมืองและอำเภออื่นภายในจังหวัดขอนแก่น

การประเมินผลของระบบมุ่งเน้นด้าน ความถูกต้องและความเหมาะสมของแผนการเดินทาง ความถูกต้องของการตอบคำถามของระบบแชตบอต ประสิทธิภาพของการค้นคืนข้อมูลผ่าน RAG ระยะเวลาในการสร้างแผนการเดินทาง และความพึงพอใจของผู้ใช้งาน โดยประเมินผ่านการทดสอบเชิงประสิทธิภาพและแบบสอบถามจากผู้ใช้งานจริง

\subsection{ข้อจำกัดของการวิจัย (Limitations)}

\begin{enumerate}[leftmargin=1.8cm]
\item ข้อจำกัดด้านพื้นที่และข้อมูล การวิจัยจำกัดขอบเขตข้อมูลเฉพาะจังหวัดขอนแก่น ส่งผลให้การนำระบบไปประยุกต์ใช้กับจังหวัดอื่นจำเป็นต้องรวบรวมและจัดเตรียมฐานข้อมูลใหม่ นอกจากนี้ ความครบถ้วนและความถูกต้องของข้อมูลขึ้นอยู่กับแหล่งข้อมูลภายนอกและข้อมูลที่ผู้ดูแลระบบนำเข้าสู่ระบบ
\item ข้อจำกัดด้านการอัปเดตข้อมูล ข้อมูลสถานที่ท่องเที่ยว ร้านอาหาร และกิจกรรมอาจมีการเปลี่ยนแปลงอยู่ตลอดเวลา เช่น เวลาเปิด--ปิด การย้ายสถานที่ หรือการยกเลิกกิจกรรม หากไม่มีการปรับปรุงข้อมูลอย่างสม่ำเสมอ อาจส่งผลให้คำแนะนำของระบบไม่สอดคล้องกับสถานการณ์จริง
\item ข้อจำกัดด้านเทคนิคของโมเดลปัญญาประดิษฐ์ แม้ว่าระบบจะใช้สถาปัตยกรรม RAG เพื่อลดปัญหาการสร้างข้อมูลที่คลาดเคลื่อน (Hallucination) ของ LLMs แต่ยังคงมีโอกาสที่โมเดลจะสร้างคำตอบที่ไม่ถูกต้องหรือไม่สอดคล้องกับบริบทของข้อมูลได้ โดยเฉพาะในกรณีที่ฐานความรู้ไม่มีข้อมูลรองรับหรือข้อมูลไม่เพียงพอ
\item ข้อจำกัดด้านค่าใช้จ่ายในการดำเนินงาน ระบบมีต้นทุนในการให้บริการ เช่น ค่าใช้จ่ายของเซิร์ฟเวอร์ ค่าใช้บริการโมเดล LLM และค่าใช้จ่ายในการจัดเก็บฐานข้อมูลเวกเตอร์ ซึ่งอาจส่งผลต่อความสามารถในการขยายระบบหรือรองรับผู้ใช้งานจำนวนมากในอนาคต
\item ข้อจำกัดของระบบสะสมคะแนน (Point System) ระบบสะสมคะแนนจากการสแกน QR Code เป็นความสามารถเพิ่มเติมของระบบ ซึ่งประสิทธิภาพในการสร้างการมีส่วนร่วมของผู้ใช้งานขึ้นอยู่กับจำนวนสถานที่ที่เข้าร่วมโครงการ ความร่วมมือจากหน่วยงานที่เกี่ยวข้อง และการดำเนินกิจกรรมส่งเสริมการท่องเที่ยวของจังหวัดขอนแก่น
\end{enumerate}
